\SourceLeaf{336}%
straight lines passes through [the] point whose polar coordinates are $\frac{1}{(12)}$, 12 (namely distance $\frac{1}{(12)}$, direction 12) and there forms with (12) the angle $02-01$, or has azimuth $21-(02-01)$.

The construction is most elegant if one puts $(12)\times(13)$ in place of 1 in the numerators of the fractions.

\medskip

\Needspace{13\baselineskip}\begin{center}[7.]\end{center}

Given points 1, 2, 3, with angles $A, B, C$ at them and sides $23=a$, $31=b$, $12=c$. The angles $\mathfrak{A}, \mathfrak{B}, \mathfrak{C}$ are prescribed at a fourth point.

$\beta$ is the circle consisting of the segments

$\beta'$, where chord $b$ is seen at angle $\mathfrak{B}$,

$\beta''$, where that chord is seen at angle $\mathfrak{B}\pm180^0$.

$\gamma$ is the circle consisting of the segments

$\gamma'$, where chord $c$ is seen at angle $\mathfrak{C}$,

$\gamma''$, where chord $c$ is seen at angle $\mathfrak{C}\pm180^0$.

The circles intersect at $A$ and $D$.

At $D$, let the angles subtended by $b$ and $c$ be $\mathfrak{B}^*, \mathfrak{C}^*$. The sines of $\mathfrak{B}^*-B$, $\mathfrak{C}^*-C$, $\mathfrak{B}^*+\mathfrak{C}^*-B-C$ will all be positive if I) $D$ lies inside the circle through 1, 2, 3, and all negative if II) $D$ lies outside it.

\Needspace{9\baselineskip}There are four possible cases for the position of $D$ in the circles $\beta$ and $\gamma$. The point $D$ lies

1) in $\beta'$ and $\gamma'$, in which case $\mathfrak{B}^*=\mathfrak{B}$, $\mathfrak{C}^*=\mathfrak{C}$

2) in $\beta''$ and $\gamma'$, where $\mathfrak{B}^*=\mathfrak{B}\pm180^0$, $\mathfrak{C}^*=\mathfrak{C}$

3) in $\beta'$ and $\gamma''$, where $\mathfrak{B}^*=\mathfrak{B}$, $\mathfrak{C}^*=\mathfrak{C}\pm180^0$

4) in $\beta''$ and $\gamma''$, where $\mathfrak{B}^*=\mathfrak{B}\pm180^0$, $\mathfrak{C}^*=\mathfrak{C}\pm180^0$.

In case 1) the problem is physically possible, and in the other three it is impossible.
